\documentclass[11pt]{article}
\def\figlang{fr}
\usepackage[a4paper,top=2.4cm,bottom=2.4cm,left=2.3cm,right=2.3cm,headsep=0.8cm,footskip=1cm]{geometry}
% Compatible XeLaTeX / LuaLaTeX (recommandé) et pdfLaTeX
\usepackage{iftex}
\ifPDFTeX
  \usepackage[T1]{fontenc}
  \usepackage[utf8]{inputenc}
  \usepackage{textcomp}
  \IfFileExists{tgtermes.sty}{\usepackage{tgtermes}}{\usepackage{mathptmx}}
  \IfFileExists{tgheros.sty}{\usepackage{tgheros}}{\usepackage[scaled=0.92]{helvet}}
\else
  \usepackage{fontspec}
  \setmainfont{TeX Gyre Termes}
  \setsansfont{TeX Gyre Heros}
\fi
\usepackage{amsmath,amssymb}
\usepackage{graphicx,booktabs,tabularx,array,multirow,float,enumitem}
\usepackage[table]{xcolor}
% Version de soumission : texte en noir sur fond blanc
\definecolor{navy}{HTML}{000000}\definecolor{ora}{HTML}{000000}\definecolor{fond}{HTML}{FFFFFF}
\usepackage[font=small,labelfont={bf,sf,color=navy},labelsep=period,justification=justified,singlelinecheck=false]{caption}
\captionsetup[table]{position=top,skip=4pt}
\usepackage{titlesec}
\titleformat{\section}{\sffamily\large\bfseries\color{navy}}{\thesection.}{0.5em}{}
\titleformat{\subsection}{\sffamily\normalsize\bfseries\color{navy}}{\thesubsection.}{0.5em}{}
\titleformat{\subsubsection}{\sffamily\normalsize\itshape\color{navy}}{\thesubsubsection.}{0.5em}{}
\titlespacing*{\section}{0pt}{14pt}{6pt}\titlespacing*{\subsection}{0pt}{10pt}{4pt}\titlespacing*{\subsubsection}{0pt}{8pt}{3pt}
\usepackage{fancyhdr}\pagestyle{fancy}\fancyhf{}
\renewcommand{\headrulewidth}{0.4pt}
\usepackage{tcolorbox}
\usepackage[hidelinks]{hyperref}
\usepackage{microtype}
\setlength{\parskip}{2pt}\setlength{\emergencystretch}{2em}\setlength{\parindent}{1.2em}
\linespread{1.06}
\setlist{itemsep=1pt,topsep=3pt,leftmargin=1.6em}
\renewcommand{\arraystretch}{1.15}
\newcommand{\fig}[3]{\begin{figure}[!htbp]\centering\includegraphics[width=#2\linewidth]{figures/\figlang_#1.png}\caption{#3}\label{fig:#1}\end{figure}}
\newcolumntype{R}{>{\raggedleft\arraybackslash}X}
\newcolumntype{L}{>{\raggedright\arraybackslash}X}

\ifPDFTeX\usepackage[french]{babel}\else\usepackage{polyglossia}\setmainlanguage{french}\fi
\fancyhead[L]{\sffamily\footnotesize\color{navy} Pluie et fièvre typhoïde dans l'Extrême-Nord du Cameroun}
\fancyhead[R]{\sffamily\footnotesize\color{navy} Matériel supplémentaire S1}
\fancyfoot[C]{\sffamily\footnotesize\thepage}
\hypersetup{colorlinks=true,linkcolor=navy,urlcolor=navy}\renewcommand{\theequation}{S\arabic{equation}}
\begin{document}\sloppy
\thispagestyle{fancy}
\noindent{\sffamily\footnotesize\color{ora}\textbf{MATÉRIEL SUPPLÉMENTAIRE S1}}\par\vspace{2pt}
\noindent{\color{navy}\rule{\linewidth}{1.2pt}}\par\vspace{8pt}
{\raggedright\sffamily\bfseries\LARGE\color{navy}Matériel supplémentaire S1\par}\vspace{4pt}
{\raggedright\sffamily\large\color{navy}Formulation mathématique complète\par}\vspace{10pt}
\noindent Ce document rassemble les formules de l'article « Influence de la pluviométrie sur la dynamique et la prévision de la fièvre typhoïde dans la région de l'Extrême-Nord du Cameroun ». Les équations sont numérotées avec le préfixe S. Les tableaux, figures et références sont présentés dans l'article principal ; les numéros entre crochets renvoient à sa liste de références.
\section{Prétraitement des données}


\subsection{Extraction et agrégation}
Soit $x_{f,k,t}$ le nombre de cas déclarés par la formation sanitaire $f$, dans la catégorie d'âge et de sexe $k$, la semaine $t$. Les 522 formations ont été rattachées à leur district $d$ grâce à la hiérarchie des unités d'organisation, et les cas agrégés :
\begin{equation}
y_{d,t}=\sum_{f\in d}\;\sum_{k} x_{f,k,t}, \qquad Y_t=\sum_{d=1}^{32} y_{d,t}.
\label{eq:agg}
\end{equation}
Une semaine-district est considérée comme manquante lorsqu'aucune formation du district n'a déclaré de valeur.

\subsection{Contrôle de qualité et harmonisation}
L'export a été comparé à une extraction antérieure portant sur 2019-2023 : les totaux régionaux hebdomadaires coïncident (écart médian de 0\,\%), les rares écarts correspondant à des corrections tardives. Cinq districts créés récemment (Gawar, Kaï-Kaï, Mokong, Tourou, Waza) ont été ramenés à leur district d'origine. Les deux formations aujourd'hui rattachées à Waza ont été réintégrées dans Mora : la somme Mora + Waza reproduit la série historique de Mora pour 250 semaines sur 261. Le nombre $N_t$ de formations déclarantes a été calculé pour chaque semaine.

\subsection{Données manquantes}
La structure des données manquantes a été examinée par district et par année. Les semaines manquantes ont été imputées par interpolation linéaire au sein de chaque district : pour une semaine manquante $t$ encadrée par les semaines observées $t_1<t<t_2$,
\begin{equation}
\hat{y}_{d,t}=y_{d,t_1}+\frac{t-t_1}{t_2-t_1}\left(y_{d,t_2}-y_{d,t_1}\right),
\label{eq:interp}
\end{equation}
puis arrondie à l'entier. Les semaines imputées ont servi à construire les variables explicatives, mais jamais de valeur de référence pour l'évaluation des prévisions par district.

\subsection{Valeurs aberrantes}
Les valeurs aberrantes ont été recherchées par le filtre de Hampel [12] : une observation est signalée lorsqu'elle s'écarte de la médiane mobile $\tilde{m}_{d,t}$ (fenêtre de 9 semaines) de plus de trois fois l'écart absolu médian normalisé,
\begin{equation}
\left|y_{d,t}-\tilde{m}_{d,t}\right|>3\times1{,}4826\times\operatorname{méd}_{|j|\le4}\left|y_{d,t+j}-\tilde{m}_{d,t}\right|.
\label{eq:hampel}
\end{equation}
Les valeurs signalées n'ont pas été corrigées : en l'absence d'information externe permettant de distinguer une erreur de saisie d'une flambée réelle, nous avons préféré les conserver, la MAE et les modèles à base d'arbres étant peu sensibles aux valeurs isolées.

\subsection{Alignement des données climatiques}
Les valeurs quotidiennes ont été agrégées en semaines du lundi au dimanche, alignées sur les semaines épidémiologiques. Pour la pluie, le cumul hebdomadaire est $P_{d,t}=\sum_{j\in t}p_{d,j}$ ; pour la température et l'humidité, la moyenne. Les séries régionales sont la moyenne des 32 districts.

\subsection{Relation moyenne-variance et transformation}
La relation entre la moyenne $\mu_d$ et la variance $\sigma^2_d$ des cas des districts a été estimée par la loi de puissance de Taylor [13] :
\begin{equation}
\log\sigma^2_d=\log a+b\,\log\mu_d .
\label{eq:taylor}
\end{equation}
Une valeur $b>1$ indique une surdispersion croissante avec le niveau des cas, ce qui justifie une transformation logarithmique $u_{d,t}=\ln(1+y_{d,t})$ lorsque des districts de tailles très différentes sont modélisés ensemble.

\section{Analyse exploratoire des séries}
\subsection{Décomposition}
La série régionale a été décomposée par la méthode STL [14] (décomposition saisonnière par régression locale, période de 52 semaines, version robuste) en une tendance $T_t$, une composante saisonnière $S_t$ et un résidu $R_t$ :
\begin{equation}
Y_t=T_t+S_t+R_t .
\label{eq:stl}
\end{equation}
La force de la tendance et celle de la saisonnalité ont été mesurées par [23] :
\begin{equation}
F_T=\max\!\left(0,\,1-\frac{\operatorname{Var}(R_t)}{\operatorname{Var}(T_t+R_t)}\right),\qquad F_S=\max\!\left(0,\,1-\frac{\operatorname{Var}(R_t)}{\operatorname{Var}(S_t+R_t)}\right).
\label{eq:force}
\end{equation}

\subsection{Indices saisonniers}
Pour quantifier la saisonnalité à l'échelle du mois, les cas ont été rapportés à leur moyenne mobile centrée sur 52 semaines, puis moyennés par mois $m$ et normalisés à une moyenne de 1, selon une démarche proche de celle de Shashvat et al. [7] :
\begin{equation}
\tilde{Y}_t=\frac{Y_t}{\frac{1}{52}\sum_{j=-26}^{25}Y_{t+j}},\qquad I_m=\frac{\overline{\tilde{Y}}_{m}}{\frac{1}{12}\sum_{m'=1}^{12}\overline{\tilde{Y}}_{m'}} .
\label{eq:indice}
\end{equation}
Un indice $I_m>1$ signale un mois au-dessus de la tendance.

\subsection{Stationnarité}
La stationnarité a été évaluée par deux tests aux hypothèses nulles opposées. Le test de Dickey-Fuller augmenté (ADF) [15] estime la régression
\begin{equation}
\Delta Y_t=\alpha+\gamma\,Y_{t-1}+\sum_{i=1}^{p}\delta_i\,\Delta Y_{t-i}+\varepsilon_t ,
\label{eq:adf}
\end{equation}
et teste $H_0:\gamma=0$ (racine unitaire, non-stationnarité). Le test KPSS [16] décompose la série en une marche aléatoire et un terme stationnaire, et teste $H_0$ : stationnarité, à l'aide de la statistique
\begin{equation}
\eta=\frac{1}{n^{2}\,\hat{\sigma}^{2}_{\infty}}\sum_{t=1}^{n}\Big(\sum_{s=1}^{t}\hat{e}_s\Big)^{2},
\label{eq:kpss}
\end{equation}
où $\hat{e}_s$ sont les résidus de la régression sur une constante et $\hat{\sigma}^{2}_{\infty}$ la variance de long terme. Une série est jugée stationnaire lorsque l'ADF rejette $H_0$ et que le KPSS ne la rejette pas. Les tests ont été appliqués à la série en niveau, en logarithme, en différence première et en différence saisonnière, ainsi qu'à chaque district.

\subsection{Autocorrélation}
Les fonctions d'autocorrélation (ACF) et d'autocorrélation partielle (PACF) ont été examinées jusqu'à 60 et 30 semaines. L'absence d'autocorrélation (des données ou des résidus d'un modèle) a été testée par la statistique de Ljung-Box [17] :
\begin{equation}
Q(m)=n(n+2)\sum_{k=1}^{m}\frac{\hat{\rho}_k^{2}}{n-k}\;\sim\;\chi^{2}_{m} .
\label{eq:lb}
\end{equation}

\section{Cadre mathématique}
\subsection{Point de départ : un modèle de transmission à réservoir environnemental}
Nous partons du modèle de Tilahun et al. [5], repris par Sabir et al. [6]. La population est divisée en sensibles $S$, porteurs asymptomatiques $C$, malades $I$ et guéris $R$ ; $B$ est la concentration de bactéries dans l'eau et les aliments :
\begin{equation}
\begin{aligned}
\frac{dS}{dt}&=\Lambda-(\lambda+\mu)S+\delta R, &\qquad \frac{dC}{dt}&=\rho\lambda S-(\phi+\theta+\mu+\sigma_1)C,\\
\frac{dI}{dt}&=(1-\rho)\lambda S+\theta C-(\alpha+\sigma_2+\mu+\beta)I, & \frac{dR}{dt}&=\phi C+\beta I-(\mu+\delta)R,\\
\frac{dB}{dt}&=\sigma_1 C+\sigma_2 I-\mu_b B .
\end{aligned}
\label{eq:seir}
\end{equation}
$\Lambda$ est le recrutement, $\mu$ la mortalité naturelle, $\rho$ la proportion des infections qui deviennent des portages, $\theta$ le passage du portage à la maladie, $\phi$ et $\beta$ les taux de guérison, $\delta$ la perte d'immunité, $\alpha$ la mortalité due à la maladie, $\sigma_1$ et $\sigma_2$ l'excrétion de bactéries et $\mu_b$ leur disparition. La force d'infection est saturante :
\begin{equation}
\lambda(t)=\frac{\nu\,B(t)}{K+B(t)} ,
\label{eq:lambda}
\end{equation}
où $\nu$ est le taux de contact maximal avec l'eau ou les aliments contaminés et $K$ la concentration de demi-saturation.

\subsection{Introduction de la pluie}
Le modèle~\eqref{eq:seir} ne contient pas de climat. Nous faisons l'hypothèse de travail que la pluie $P(t)$ augmente la contamination de l'environnement, par ruissellement et contamination des puits et des points d'eau, avec une intensité $\kappa\ge0$ :
\begin{equation}
\frac{dB}{dt}=\sigma_1 C+\sigma_2 I+\kappa P(t)-\mu_b B,
\qquad
B(t)=\int_0^{\infty}e^{-\mu_b s}\big[\sigma_1C(t-s)+\sigma_2I(t-s)+\kappa P(t-s)\big]\,ds .
\label{eq:pluie}
\end{equation}
La concentration en bactéries intègre ainsi toute la pluie passée, avec une mémoire de durée moyenne $1/\mu_b$.

\subsection{Des infections aux cas déclarés}
Une infection ne devient un cas déclaré qu'après l'incubation (environ 6 à 30 jours [6]), la consultation et la déclaration. Avec $g(s)$ la distribution de ce délai total et $\pi_t$ la proportion des cas captés par la surveillance, qui croît avec le nombre $N_t$ de formations déclarantes :
\begin{equation}
\mathbb{E}[y_t]=\pi_t\int_0^{\infty}g(s)\big[(1-\rho)\lambda(t-s)S(t-s)+\theta C(t-s)\big]\,ds,\qquad \pi_t=\pi(N_t),\;\pi'(N)>0 .
\label{eq:obs}
\end{equation}

\subsection{Prédictions qualitatives}
Les équations~\eqref{eq:lambda} à~\eqref{eq:obs} conduisent à trois prédictions. (i) \textit{Un décalage entre pluie et cas} : la pluie agit à travers la mémoire du réservoir ($1/\mu_b$) puis le délai $g$, si bien que le pic des cas suit celui des pluies de plusieurs semaines. (ii) \textit{Un effet saturant} : la sensibilité $\partial\lambda/\partial B=\nu K/(K+B)^2$ tend vers zéro quand $B\gg K$, de sorte qu'un surplus de pluie peut ne plus accroître le risque pendant la saison des pluies. (iii) \textit{Des tendances liées à la surveillance} : une hausse de $\pi_t$ accroît les cas déclarés sans changement de l'incidence réelle. Les paramètres ($\nu$, $K$, $\kappa$, $\mu_b$, $g$) ne sont pas connus localement et $S$, $C$, $B$ ne sont pas observés : le modèle n'est donc pas ajusté, mais sert à dériver les hypothèses testées statistiquement.

\section{Formulation du problème de prévision et hypothèses}\label{sec:hyp}
Soit $\mathcal{F}_t$ l'information disponible à la fin de la semaine $t$. La précision étant mesurée par l'erreur absolue, la prévision optimale à $h$ semaines est la médiane conditionnelle :
\begin{equation}
f_h^{*}=\arg\min_f\;\mathbb{E}\big|y_{t+h}-f(\mathcal{F}_t)\big|=\operatorname{méd}\big(y_{t+h}\mid\mathcal{F}_t\big).
\label{eq:objectif}
\end{equation}
Avec $\tilde{Y}_t$ défini en~\eqref{eq:indice}, $P^{(4)}_t=\sum_{j=0}^{3}P_{t-j}$ le cumul de pluie sur 4 semaines et $A_t=P^{(4)}_t-\overline{P^{(4)}}_{w(t)}$ son anomalie par rapport à la même semaine calendaire $w(t)$, les hypothèses s'écrivent :
\begin{align}
\textbf{H1 :}&\quad r(\ell)=\operatorname{corr}\big(\tilde{Y}_t,P^{(4)}_{t-\ell}\big)>0,\ \text{maximale pour un décalage }\ell^{*}>0 ; \label{eq:h1}\\
\textbf{H2 :}&\quad \operatorname{corr}\big(\tilde{Y}_t,A_{t-\ell^{*}}\big)>0 ; \label{eq:h2}\\
\textbf{H3 :}&\quad \mathbb{E}\big|y_{t+h}-\hat{f}_h(\mathcal{F}_t,\mathbf{z}_t)\big|<\mathbb{E}\big|y_{t+h}-\hat{f}_h(\mathcal{F}_t)\big| ; \label{eq:h3}\\
\textbf{H4 :}&\quad \mathbb{E}\big|e^{\text{global}}_{d,t+h}\big|<\mathbb{E}\big|e^{\text{naïf}}_{d,t+h}\big|\ \text{et}\ \mathbb{E}\big|e^{\text{global}}_{d,t+h}\big|<\mathbb{E}\big|e^{\text{local}}_{d,t+h}\big| , \label{eq:h4}
\end{align}
où $\mathbf{z}_t$ désigne les variables climatiques observées jusqu'à la semaine $t$. Si la médiane conditionnelle de $y_{t+h}$ ne change pas lorsqu'on ajoute $\mathbf{z}_t$, toute l'information utile du climat est déjà contenue dans $\mathcal{F}_t$ et H3 est rejetée.

\section{Modèles de prévision}
\textbf{Références naïves.}
\begin{equation}
\hat{y}^{\,\text{naïf}}_{t+h|t}=y_t,\qquad \hat{y}^{\,\text{saison}}_{t+h|t}=y_{t+h-52}.
\label{eq:naif}
\end{equation}
\textbf{ARIMA} [18]. Avec $L$ l'opérateur de retard et $\varepsilon_t$ un bruit blanc,
\begin{equation}
\Big(1-\sum_{i=1}^{p}\phi_iL^{i}\Big)(1-L)^{d}\,y_t=c+\Big(1+\sum_{j=1}^{q}\theta_jL^{j}\Big)\varepsilon_t .
\label{eq:arima}
\end{equation}
L'ordre $d$ a été choisi par le test KPSS, $p$ et $q$ (0 à 3) par le critère d'Akaike sur la période d'entraînement initiale, puis figés.

\textbf{Prophet} [19]. Tendance linéaire par morceaux $g(t)$ et saisonnalité annuelle de Fourier d'ordre 10, de période $P=52{,}18$ semaines :
\begin{equation}
y(t)=g(t)+\sum_{k=1}^{10}\Big[a_k\cos\frac{2\pi kt}{P}+b_k\sin\frac{2\pi kt}{P}\Big]+\varepsilon_t .
\label{eq:prophet}
\end{equation}
\textbf{XGBoost} [20]. La prévision est la somme de $M$ arbres de régression construits séquentiellement ; chaque arbre minimise une perte quadratique pénalisée par sa complexité ($T$ feuilles de valeurs $w$) :
\begin{equation}
\hat{y}=\sum_{m=1}^{M}f_m(\mathbf{x}),\qquad \mathcal{L}=\sum_i(y_i-\hat{y}_i)^2+\sum_{m=1}^{M}\Big(\gamma T_m+\tfrac{1}{2}\lambda_r\lVert w_m\rVert^2\Big).
\label{eq:xgb}
\end{equation}
Compte tenu de la non-stationnarité en niveau (voir l'article principal), nous avons retenu une stratégie directe qui prédit la \textit{variation} de la série, avec un modèle $F_h$ par horizon :
\begin{equation}
\hat{y}_{t+h|t}=y_t+F_h(\mathbf{x}_t),\qquad \mathbf{x}_t=\Big(y_t,\dots,y_{t-3},\ \bar{y}^{(4)}_t,\ \bar{y}^{(13)}_t,\ \sin\tfrac{2\pi w_{t+h}}{52},\ \cos\tfrac{2\pi w_{t+h}}{52},\ y_{t+h-52}\Big).
\label{eq:direct}
\end{equation}
Le modèle \textbf{XGBoost + climat} complète $\mathbf{x}_t$ par $\mathbf{z}_t=\big(P^{(4)}_t,\,P^{(4)}_{t-4},\,\bar{T}^{(4)}_t,\,\bar{H}^{(4)}_t\big)$, toutes observées à la date de prévision.

\textbf{Modèles local et global.} Le modèle local ajuste une fonction $F_{h,d}$ par district. Le modèle global ajuste une seule fonction $F_h$ sur les 32 districts, en échelle logarithmique, avec le niveau moyen passé du district $\bar{u}_{d,t}$ :
\begin{equation}
u_{d,t+h}-u_{d,t}=F_h\big(\mathbf{x}_{d,t},\bar{u}_{d,t}\big)+\varepsilon_{d,t+h},\qquad \hat{y}_{d,t+h|t}=\exp\big(u_{d,t}+\hat{F}_h\big)-1 .
\label{eq:global}
\end{equation}
Les hyperparamètres ont été fixés a priori et jamais ajustés sur la période test.



\section{Protocole d'évaluation}
Les prévisions ont été évaluées hors échantillon par origine glissante, avec fenêtre d'entraînement croissante [23]. La période test couvre 2023-2025 (156 semaines), choisie parce que le réseau de surveillance y est stable. Pour chaque semaine $t$ de cette période, chaque modèle est réentraîné sur les seules semaines $1,\dots,t$, puis prévoit les cas à $h=1$ et $h=4$ semaines. Avec $e_{t+h}=y_{t+h}-\hat{y}_{t+h|t}$ :
\begin{equation}
\text{MAE}=\frac{1}{n}\sum|e_{t+h}|,\quad \text{RMSE}=\sqrt{\frac{1}{n}\sum e_{t+h}^2},\quad \text{MAPE}=\frac{100}{n}\sum\frac{|e_{t+h}|}{y_{t+h}},\quad \text{MAE}^{\text{rel}}_d=\frac{\text{MAE}^{\text{modèle}}_d}{\text{MAE}^{\text{naïf}}_d}.
\label{eq:metrics}
\end{equation}
La MAE est le critère principal, car elle est moins sensible aux valeurs extrêmes que la RMSE [24].

\section{Tests statistiques}
Pour comparer deux modèles A et B, on calcule $d_t=|e^{A}_t|-|e^{B}_t|$, sa moyenne $\bar{d}$ et ses autocovariances $\gamma_k$. La statistique de Diebold-Mariano [25], corrigée pour les petits échantillons selon Harvey, Leybourne et Newbold [26], est comparée à une loi de Student à $n-1$ degrés de liberté :
\begin{equation}
\mathrm{DM}=\frac{\bar{d}}{\sqrt{\frac{1}{n}\big(\gamma_0+2\sum_{k=1}^{h-1}\gamma_k\big)}},\qquad \mathrm{DM}^{*}=\mathrm{DM}\sqrt{\frac{n+1-2h+h(h-1)/n}{n}}\;\sim\;t_{n-1}.
\label{eq:dm}
\end{equation}
Pour l'ensemble des districts, l'écart d'erreur est normalisé par la MAE naïve de chaque district puis moyenné chaque semaine sur les $D_t$ districts observés ; le test~\eqref{eq:dm} est appliqué à cette série :
\begin{equation}
\tilde{d}_t=\frac{1}{D_t}\sum_{d}\frac{|e^{A}_{d,t}|-|e^{B}_{d,t}|}{\text{MAE}^{\text{naïf}}_d}.
\label{eq:panel}
\end{equation}
Les corrélations~\eqref{eq:h1} et~\eqref{eq:h2} sont des corrélations de Pearson ; les associations entre caractéristiques des districts ont été testées par corrélation de Spearman. Le seuil de significativité est de 5\,\%.


\end{document}
